\section*{Chapter \ref{ch:g10:market-equilibrium} --- Market Equilibrium}
% figdata: figdata/grade-10/06-market-equilibrium/out/ (world_oil, world_moves); Brenland and Kestria numbers invented.

\begin{solution}{exo:g10:market-equilibrium:1}
$120 - 4p = 6p - 30$ gives $10p = 150$, $p^* = 15$, $Q^* = 120 - 60 = 60$ (check: $90 - 30 = 60$).
\end{solution}

\begin{solution}{exo:g10:market-equilibrium:2}
$90 - p = 2p$ gives $p^* = 30$, $Q^* = 60$. Check: $90 - 30 = 60$ and $2 \times 30 = 60$.
\end{solution}

\begin{solution}{exo:g10:market-equilibrium:3}
At 8 dalens: $Q^d = 26$, $Q^s = 6$, an \omterm{def:g10:market-equilibrium:excess-demand}{excess demand} of 20. At 11: $Q^d = 17$, $Q^s = 12$, an \omterm{def:g10:market-equilibrium:excess-demand}{excess demand} of 5. At 13: $Q^d = 11$, $Q^s =
16$, an \omterm{def:g10:market-equilibrium:excess-demand}{excess supply} of 5. (Each dalen changes the gap by $b + d = 5$, as \cref{prop:g10:market-equilibrium:excess} says.)
\end{solution}

\begin{solution}{exo:g10:market-equilibrium:4}
In 2012 (\usd{4.98}) and 2011 (\usd{4.92}). Between 2016 and 2018 both rose: the price from \usd{2.80} to \usd{3.42} and production from 97.09 to
100.88 million barrels a day.
\end{solution}

\begin{solution}{exo:g10:market-equilibrium:5}
(a) Supply, right. (b) Demand, right: more people want to go to the coast. (c) Supply, left. (d) Demand, right: the town is easier to reach.
\end{solution}

\begin{solution}{exo:g10:market-equilibrium:6}
$50 - 3p = 2p - 20$ gives $5p = 70$, $p^* = 14$, $Q^* = 50 - 42 = 8$. By \cref{prop:g10:market-equilibrium:one-shift} with $\Sigma = -10$: the price
changes by $10/5 = 2$ (from 12 to 14) and the quantity by $3 \times (-10)/5 = -6$ (from 14 to 8).
\end{solution}

\begin{solution}{exo:g10:market-equilibrium:7}
Price: $(15 - 5)/5 = 2$, so $p^* = 14$. Quantity: $(2 \times 15 + 3 \times 5)/5 = 45/5 = 9$, so $Q^* = 23$. Check: $65 - 3p = 2p - 5$ gives $5p = 70$,
$p = 14$, $Q = 65 - 42 = 23 = 28 - 5$.
\end{solution}

\begin{solution}{exo:g10:market-equilibrium:8}
Coffee: supply shifts left; its price rises and the quantity falls. Tea: dearer coffee shifts the demand for tea right; its price and quantity both
rise. (Drawn as the two panels of the figure on apples and pears.)
\end{solution}

\begin{solution}{exo:g10:market-equilibrium:9}
United States: $15.15 - 8.56 = 6.59$ million barrels a day. Rest of the world: $81.50 - 78.34 = 3.16$. The world: $6.59 + 3.16 = 9.75$, so the
United States' share was $6.59/9.75 \approx 68\,\%$, about two-thirds.
\end{solution}

\begin{solution}{exo:g10:market-equilibrium:10}
The \omterm{def:g10:supply:market-supply}{market supply} is $p$ for $p < 4$, $2p - 4$ for $4 \le p < 10$ and $4p - 24$ for $p \ge 10$. On the middle piece, $48 - 2p = 2p - 4$ gives $p = 13$,
which is not below 10: the equilibrium is not there. On the top piece, $48 - 2p = 4p - 24$ gives $6p = 72$, $p^* = 12$, which is above 10, and $Q^* =
48 - 24 = 24$. Check: $12 + 8 + 4 = 24$.
\end{solution}

\begin{solution}{exo:g10:market-equilibrium:11}
At the old price of 12, demand is $60 - 36 = 24$ and supply 14: an \omterm{def:g10:market-equilibrium:excess-demand}{excess demand} of 10. The price rises by 2 dalens. Buyers move along the new demand
curve and want $3 \times 2 = 6$ fewer (from 24 to 18); sellers move along the supply curve and offer $2 \times 2 = 4$ more (from 14 to 18). The gap of
10 is closed by $6 + 4$.
\end{solution}

\begin{solution}{exo:g10:market-equilibrium:12}
A rise of both the price and the quantity shows that demand shifted right; it does not show that supply stayed put. By
\cref{prop:g10:market-equilibrium:two-shifts}, supply may also have shifted right, by less than demand (the price still rises), or even left, by
little enough for the quantity still to rise. The observation tells the direction of the demand shift and that it dominated; it says nothing
certain about supply.
\end{solution}

\begin{solution}{exo:g10:market-equilibrium:13}
By \cref{prop:g10:market-equilibrium:one-shift}, the demand shift changes the price by $-\Delta/(b+d)$ and the supply shift by $-\Sigma/(b+d)$: the
total $-(\Delta + \Sigma)/(b + d)$ is negative, so the price falls. The quantity changes by $(-d\Delta + b\Sigma)/(b+d)$, which is positive when $b\Sigma >
d\Delta$: the supply shift, weighted by the slope of demand, outweighs the demand shift, weighted by the slope of supply.
\end{solution}

\begin{solution}{exo:g10:market-equilibrium:14}
At 10 dalens bakers offer $2 \times 10 - 10 = 10$ hundred loaves and buyers want $50 - 30 = 20$: only 10 hundred loaves are sold, the smaller of the
two. The first customers get them; half the bread wanted at that price is not baked. Queues before opening and empty shelves by mid-morning are what
one would expect to see.
\end{solution}

\begin{solution}{exo:g10:market-equilibrium:15}
$a - bp = c - dp$ gives $(a - c) = (b - d)\,p$. If $d = b$: when $a \ne c$ there is no solution (the lines are parallel); when $a = c$ every price is a
solution (the lines coincide). If $d > b$ and $a > c$: $p = (a - c)/(b - d)$, a positive number divided by a negative one, so the price is negative and
the market has no equilibrium at a positive price. The proof divides by $b + d$ (here $b - d$) and uses that it is positive; with a supply curve
sloping down as steeply as demand or more, that step fails.
\end{solution}

\begin{solution}{pb:g10:market-equilibrium:1}
Prices are in kests per kilogram, quantities in tonnes a week.

\textbf{1.} It rises by one tonne a week: a dearer apple sends buyers to pears, the mark of \omterm{def:g10:demand:substitutes}{substitutes}.

\textbf{2.} $100 - 2p_A + p_P = 3p_A - 20$, that is $5p_A - p_P = 120$.

\textbf{3.} $80 - 2p_P + p_A = 2p_P - 10$, that is $-p_A + 4p_P = 90$.

\textbf{4.} From question 2, $p_P = 5p_A - 120$. Then $-p_A + 20p_A - 480 = 90$, so $19p_A = 570$, $p_A = 30$ and $p_P = 150 - 120 = 30$.

\textbf{5.} Apples: $3 \times 30 - 20 = 70$ (demand $100 - 60 + 30 = 70$). Pears: $2 \times 30 - 10 = 50$ (demand $80 - 60 + 30 = 50$).

\textbf{6.} By $58 - 20 = 38$ tonnes a week at every price.

\textbf{7.} $100 - 2p_A + 30 = 3p_A - 58$ gives $5p_A = 188$, $p_A = 37.6$, and $Q_A = 3 \times 37.6 - 58 = 112.8 - 58 = 54.8$.

\textbf{8.} Pear demand $80 - 60 + 37.6 = 57.6$, pear supply 50: an \omterm{def:g10:market-equilibrium:excess-demand}{excess demand} of 7.6 tonnes.

\textbf{9.} Apples: $100 - 2p_A + p_P = 3p_A - 58$, that is $5p_A - p_P = 158$. Pears: $-p_A + 4p_P = 90$, unchanged.

\textbf{10.} $p_P = 5p_A - 158$, so $-p_A + 20p_A - 632 = 90$, $19p_A = 722$, $p_A = 38$ and $p_P = 190 - 158 = 32$.

\textbf{11.} Apples: supply $3 \times 38 - 58 = 56$, demand $100 - 76 + 32 = 56$. Pears: supply $2 \times 32 - 10 = 54$, demand $80 - 64 + 38 = 54$.

\textbf{12.} By 8 kests, from 30 to 38. With the pear price held at 30 it would have risen to 37.6: the pear market's response adds 0.4 kests.

\textbf{13.} Dearer apples moved the demand for pears right, by 8 tonnes at each pear price; pear growers moved along their supply curve to a higher price
and a larger quantity.

\textbf{14.} A tonne is \num{1000} kilograms. Before: $30 \times 70 \times 1000 = \kst{2100000}$ a week. After: $38 \times 56 \times 1000 = \kst{2128000}$.

\textbf{15.} Before: $30 \times 50 \times 1000 = \kst{1500000}$. After: $32 \times 54 \times 1000 = \kst{1728000}$.

\textbf{16.} The pear growers gained clearly (\kst{228000} a week more). The apple growers' takings rose a little (\kst{28000}, about 1.3\,\%): the price
rose by more, in proportion, than the quantity fell. As a whole, apple growers can gain from a frost when buyers cut their purchases little as the price
rises. It depends on how strongly demand responds to the price, and the gain is not shared: growers whose blossom froze lose, those spared gain.

\textbf{17.} Apple price $8/30 \approx +26.7\,\%$; pear price $2/30 \approx +6.7\,\%$; apples sold $-14/70 = -20\,\%$; pears sold $4/50 = +8\,\%$.

\textbf{18.} Apples: both. The frost moved the market along the apple demand curve, and the dearer pears shifted that demand right a little (by 2 tonnes at
each apple price). Pears: the demand curve shifted right, and the market moved along the pear supply curve.

\textbf{19.} The pear market would not have moved: the price of pears would have stayed at 30 (and the apple market would have stopped at 37.6 kests and 54.8
tonnes).

\textbf{20.} After the frost, apples rose from 30 to 38 kests a kilogram and the quantity sold fell from 70 to 56 tonnes a week; pears rose from 30 to 32
kests and the quantity sold rose from 50 to 54 tonnes, although no pear tree was touched.
\end{solution}
